% =========================================================================
% Capítulo 3 — Objetivos y alcance del trabajo (OBLIGATORIO)
% Objetivo principal + objetivos secundarios (medibles), qué entra y qué
% queda explícitamente FUERA del alcance (límites), y criterios de éxito.
% =========================================================================

\chapter{Objetivos y alcance del trabajo}
\label{ch:objetivos}

Los objetivos de este \ac{TFG} se formulan siguiendo el marco \ac{SMART}, es decir, como compromisos específicos, medibles, alcanzables, relevantes y acotados en el tiempo \parencite{doran-smart}.
Se opta por este marco frente a una enumeración informal de intenciones porque convierte cada objetivo en un criterio verificable: al cierre del trabajo puede comprobarse de forma objetiva si se ha cumplido o no.

El horizonte temporal común a todos los objetivos es la defensa del trabajo en la convocatoria de junio de \thesisYear.
El desarrollo se concentra en un periodo de ejecución de unos tres meses, con una dedicación de referencia de 360 horas, correspondiente a los 12 créditos \ac{ECTS} de la asignatura.
La planificación detallada de tareas, fechas e hitos se recoge en los capítulos \ref{ch:tareas} y \ref{ch:gantt-tfg}.

Los objetivos se organizan en dos bloques: los \emph{técnicos y de producto}, que describen qué debe existir y funcionar al terminar, y los \emph{de gestión y académicos}, que describen cómo debe ejecutarse y documentarse el trabajo.
Dentro del primer bloque se distingue entre objetivos comprometidos y objetivos secundarios: los segundos aportan valor y validan decisiones de diseño, pero su ausencia no compromete el éxito del trabajo.

\section{Objetivo principal}
\label{sec:objetivo-principal}

Diseñar, desarrollar y desplegar \bukit, una red social orientada a ayudar a cumplir objetivos vitales (\eng{bucket lists}), entregando antes de la defensa un \eng{backend} completo (\ac{API} \ac{REST} y base de datos) y un cliente web funcional que cubra la experiencia de producto definida, junto con una memoria que documente y justifique todas las decisiones tomadas.

\section{Objetivos técnicos y de producto}
\label{sec:objetivos-tecnicos}

\begin{itemize}
  \item \textbf{Backend único para las tres necesidades del producto.}
        Implementar con Express y Prisma una \ac{API} \ac{REST} sobre una única base de datos PostgreSQL que resuelva el grafo social (integridad relacional), el mapa (consultas geoespaciales por radio con PostGIS) y la afinidad entre personas usuarias (búsqueda por similitud vectorial con pgvector).
        El objetivo se considera cumplido cuando el 100\,\% de los flujos núcleo del producto (cuentas, listas y metas, publicaciones, seguimiento, mensajería, mapa y notificaciones) está expuesto por la \ac{API} y cubierto por pruebas de integración antes del cierre del periodo de ejecución.
  \item \textbf{Autenticación y seguridad desde el diseño.}
        Proteger la \ac{API} con autenticación \ac{JWT} con rotación de \eng{refresh tokens}, contraseñas cifradas con Bcrypt, doble capa de limitación de peticiones y manejo centralizado de errores.
        Medible: el 100\,\% de los \eng{endpoints} privados exige un token válido, y los casos de abuso básicos (credenciales inválidas, exceso de peticiones) quedan cubiertos por pruebas automatizadas.
  \item \textbf{Cliente web completo.}
        Desarrollar con React y Vite una \ac{SPA} que implemente las cuatro pestañas del producto: inicio con \eng{feed} personalizado, listas, mensajes con mapa y notificaciones, y perfil.
        Medible: una persona usuaria puede completar de principio a fin los flujos núcleo desde el navegador, sin errores bloqueantes, antes del cierre del periodo de ejecución.
  \item \textbf{Lógica compartida sin duplicación.}
        Centralizar los tipos, la validación (Zod) y el cliente HTTP en los paquetes compartidos del \eng{monorepo} (\texttt{packages/core}, \texttt{packages/api-client}, \texttt{packages/ui-core}) como única fuente de verdad para web y móvil.
        Medible: cero esquemas de validación duplicados entre aplicaciones al cierre del desarrollo.
  \item \textbf{Aplicación móvil parcial (objetivo secundario).}
        Construir con React Native y Expo las pantallas núcleo (autenticación, listas y \eng{feed}) reutilizando la lógica compartida, como demostración de la arquitectura multiplataforma.
        Se declara secundario porque el compromiso mínimo del trabajo es \eng{backend} más web; completarlo valida la decisión arquitectónica central, pero no completarlo no invalida el trabajo.
  \item \textbf{Estimación de plazos asistida por \ac{IA} (objetivo secundario).}
        Integrar la estimación de plazos realistas por meta mediante un modelo de lenguaje con salida estructurada, ejecutada en segundo plano y dirigida por eventos (nunca por sondeo).
        Se declara secundario porque depende de servicios externos cuyo coste y disponibilidad quedan fuera del control del proyecto.
  \item \textbf{Calidad y despliegue continuo.}
        Montar una cadena de integración continua con estrategia de fallo rápido (análisis estático, comprobación de tipos, pruebas unitarias y de integración, y pruebas \ac{E2E} como última barrera) con despliegue automático al superar todas las etapas.
        Medible: la rama principal se mantiene en verde y cada cambio validado llega a producción sin intervención manual.
\end{itemize}

\section{Objetivos de gestión y académicos}
\label{sec:objetivos-gestion}

\begin{itemize}
  \item \textbf{Dedicación y plazo.}
        Ejecutar el desarrollo en un periodo de unos tres meses, con una dedicación total en el entorno de las 360 horas correspondientes a los 12 créditos \ac{ECTS}, finalizando producto y memoria antes de la convocatoria de junio de \thesisYear.
        Medible: el registro de horas por sesión del diario de desarrollo del repositorio.
  \item \textbf{Trazabilidad de las decisiones.}
        Registrar el 100\,\% de las sesiones de trabajo en el diario de desarrollo, anotando fecha, horas invertidas, problemas encontrados y cada decisión junto con la alternativa descartada y su porqué.
        Este diario es la fuente de datos real de la planificación (capítulos \ref{ch:tareas} y \ref{ch:gantt-tfg}), por lo que su exhaustividad condiciona directamente la calidad de la memoria.
  \item \textbf{Memoria en paralelo al desarrollo.}
        Redactar la memoria de forma incremental, junto al desarrollo y no al final, manteniendo el documento siempre compilable sin errores.
        Medible: al cierre de cada fase del desarrollo, los capítulos que la documentan contienen redacción real y no marcadores de contenido pendiente.
  \item \textbf{Empleo documentado de la \ac{IA} generativa.}
        Emplear herramientas de \ac{IA} generativa como asistentes del desarrollo y declarar su uso conforme a la guía de la \ac{EIB}, analizando de forma crítica su aportación real, sus límites y los problemas encontrados.
        Medible: la declaración de empleo de \ac{IA} de esta memoria enumera todas las herramientas utilizadas y los usos que se les han dado.
\end{itemize}

\section{Alcance}
\label{sec:alcance}

Delimitar el alcance es tan importante como enumerar los objetivos: fija qué se compromete, asegura la entrega de un producto mínimo funcional y protege el trabajo frente a la expansión descontrolada de requisitos (\eng{feature creep}).

\subsection{Dominio de la aplicación}
\label{sec:alcance-dominio}

El sistema se contextualiza en el dominio de una red social orientada a la acción: ayudar a cumplir objetivos vitales (\eng{bucket lists}) y a no perseguirlos en solitario.
Los elementos principales sobre los que orbita el sistema son las personas usuarias, las listas de metas con plazo, las publicaciones de progreso y las conexiones entre personas afines (seguimiento, mensajería y mapa).
La experiencia se organiza en cuatro pestañas (inicio, listas, mensajes y perfil), y esa misma estructura ordena los módulos funcionales del alcance.

\subsection{Requisitos funcionales (dentro del alcance)}
\label{sec:alcance-dentro}

Para garantizar la viabilidad del trabajo, el desarrollo se limita a los siguientes módulos:

\begin{itemize}
  \item \textbf{Usuarios, roles y verificación.}
        Registro e identificación única, con dos tipos de usuario (normal y TOP) que reciben distinta ponderación en el \eng{feed}.
        La verificación de identidad se delega en un proveedor externo en modo de pruebas, almacenando únicamente el estado y el \eng{token} de la verificación por minimización de datos conforme al \ac{RGPD}; superarla desbloquea publicar desde el perfil y el mapa de descubrimiento, restringido a personas mayores de edad.
  \item \textbf{Listas y metas (unidad central del producto).}
        Creación, consulta, edición y eliminación de listas de metas con plazo, con recordatorios asociados y un máximo de tres listas en el plan gratuito.
  \item \textbf{Publicaciones y \eng{feed}.}
        Publicaciones de progreso y de meta completada, creadas únicamente desde listas propias; \eng{feed} personalizado en mosaico con ponderación por tipo de usuario; segundo \eng{feed} solo de seguidos; y buscador de perfiles, publicaciones y listas.
  \item \textbf{Mensajería, mapa y notificaciones.}
        Chat en tiempo real entre personas usuarias; mapa con vista de seguidos mutuos y vista de descubrimiento por proximidad, apoyada en el cálculo de afinidad basado en \eng{embeddings}; y centro de notificaciones.
  \item \textbf{Perfil.}
        Hasta cinco publicaciones libres e historias ilimitadas, con contenido efímero que se elimina a los 60 días salvo acción expresa de mantenerlo.
  \item \textbf{Infraestructura transversal.}
        El \eng{monorepo} completo, con la aplicación móvil en su versión parcial (objetivo secundario, sección \ref{sec:objetivos-tecnicos}) y los paquetes compartidos de lógica; la estimación de plazos asistida por \ac{IA}, dirigida por eventos; y la infraestructura de calidad: pruebas unitarias, de integración y \ac{E2E}, con integración y despliegue continuos.
\end{itemize}

\subsection{Funcionalidades excluidas (fuera del alcance)}
\label{sec:alcance-fuera}

Para acotar el esfuerzo y proteger los plazos, quedan explícitamente fuera del alcance de este \ac{TFG}, como líneas futuras:

\begin{itemize}
  \item \textbf{Planes de pago.}
        No se implementan pasarelas de pago ni lógica de suscripción; el límite de tres listas del plan gratuito queda fijo.
  \item \textbf{Verificación de identidad en producción.}
        El proveedor externo se integra únicamente en modo de pruebas; su contratación real y los costes asociados quedan fuera.
  \item \textbf{Red de anuncios real.}
        No se integra ninguna red comercial; el hueco publicitario del \eng{feed} se cubre con una tabla propia de publicaciones patrocinadas.
  \item \textbf{Soporte de vídeo.}
        No se implementan transcodificación ni \eng{streaming}; el contenido multimedia se limita a imágenes.
  \item \textbf{Publicación en tiendas de aplicaciones.}
        La distribución en App Store y Play Store queda fuera; la aplicación móvil se valida con compilaciones de desarrollo.
  \item \textbf{Funciones aún en definición.}
        La penalización social por metas incumplidas y el panel de administración de usuarios TOP se posponen hasta que su diseño madure.
\end{itemize}

Estas líneas se retoman como trabajo futuro en el capítulo \ref{ch:conclusiones}.

\section{Criterios de éxito}
\label{sec:criterios-exito}

El trabajo se considerará un éxito si, a su cierre, se cumplen de manera binaria (sí/no) las siguientes condiciones:

\begin{enumerate}
  \item \textbf{Producto funcional desplegado.}
        Una persona usuaria puede registrarse, crear una lista con metas y plazo, publicar un progreso desde su lista, seguir a otras personas, consultar su \eng{feed}, conversar por chat y usar el mapa, sin errores bloqueantes y sobre el despliegue de producción.
  \item \textbf{Arquitectura compartida demostrada.}
        Los tipos y la validación viven una única vez en los paquetes compartidos y ninguna aplicación duplica esquemas.
  \item \textbf{Calidad automatizada.}
        La cadena de integración y despliegue continuos está operativa, en verde sobre la rama principal, y cada cambio validado se despliega sin intervención manual.
  \item \textbf{Trazabilidad completa.}
        El diario de desarrollo cubre todas las sesiones de trabajo y la memoria completa el índice oficial de la \ac{EIB} sin contenidos pendientes, con las decisiones justificadas frente a sus alternativas.
  \item \textbf{Ejecución dentro del marco académico.}
        El desarrollo se ha completado dentro del periodo planificado y la dedicación registrada es coherente con las 360 horas de los 12 créditos \ac{ECTS}.
\end{enumerate}
